% ============ Appendix to §7. Pure submodules (原书 pp.53-54) ============
\section*{Appendix to \S\ref{sec:7}. Pure submodules}
\phantomsection
\addcontentsline{toc}{section}{Appendix to \S\ref{sec:7}. Pure submodules}

Let $A$ be a ring and $M$ an $A$-module. A submodule $N$ of $M$ is said to be \emph{pure}\index{pure submodule}
if the sequence $0 \to N \otimes E \to M \otimes E$ is exact for every $A$-module $E$. Since
tensor product and exactness commute with inductive limits, we need only
consider $A$-modules $E$ of finite presentation.

\begin{example}\label{ex:7.1}
If $M/N$ is a flat $A$-module, then $N$ is a pure submodule of
$M$. This follows from the exact sequence $\Tor^A_1(M/N, E) \to N \otimes E
\to M \otimes E$.
\end{example}

\begin{example}\label{ex:7.2}
Any direct summand of $M$ is a pure submodule.
\end{example}

\begin{example}\label{ex:7.3}
If $A = \Z$, a submodule $N$ of $M$ is pure if and only if $N \cap mM
= mN$ for all $m > 0$. In fact the condition is equivalent to the exactness
of $0 \to N \otimes \Z/m\Z \to M \otimes \Z/m\Z$, and every finitely generated $\Z$-module
is a direct sum of cyclic modules.
\end{example}

\begin{theorem}\label{thm:7.13}
A submodule $N$ of $M$ is pure if and only if the following
condition holds: if $x_i = \sum_{j=1}^{s} a_{ij}m_j$ (for $1 \le i \le r$), with $m_j \in M$, $x_i \in N$ and
$a_{ij} \in A$, then there exist $y_j \in N$ (for $1 \le j \le s$) such that $x_i = \sum_{j=1}^{s} a_{ij}y_j$
(for $1 \le i \le r$).
\end{theorem}

\begin{proof}
Suppose $N$ is pure in $M$. Consider the free module $A^r$ with basis
$e_1, \ldots, e_r$ and let $D$ be the submodule of $A^r$ generated by $\sum_i a_{ij}e_i$, $1 \le j \le s$. Set
$E = A^r/D$, and let $\bar{e}_i$ denote the image of $e_i$ in $E$. Then in $M \otimes E$ we have
\[
  \sum_i x_i \otimes \bar{e}_i = \sum_{i,j} a_{ij}m_j \otimes \bar{e}_i = \sum_j m_j \otimes \left(\sum_i a_{ij}\bar{e}_i\right) = 0,
\]
hence $\sum_i x_i \otimes \bar{e}_i = 0$ in $N \otimes E$ by purity. But this means that, in $N \otimes A^r$, the
element $\sum_i x_i \otimes e_i$ is of the form $\sum_j y_j \otimes \sum_i a_{ij}e_i$ for some $y_j \in N$.

Conversely, suppose the condition is satisfied. Let $E$ be an $A$-module of
finite presentation. Then we can write $E = A^r/D$ with $D$ generated by a finite
number of elements of $A^r$, say $\sum_i a_{ij}e_i$, $1 \le j \le s$. Then reversing the
preceding argument we can see that $N \otimes E \to M \otimes E$ is injective.
\end{proof}

\begin{theorem}\label{thm:7.14}
If $N$ is a pure submodule and $M/N$ is of finite presentation,
then $N$ is a direct summand of $M$.
\end{theorem}

\begin{proof}
We will prove that $0 \to N \xrightarrow{i} M \xrightarrow{p} M/N \to 0$ splits, where $i$
and $p$ are the natural maps. For this we need only construct a linear map
$f\colon M/N \to M$ such that $pf$ is the identity map of $M/N$. Let $\{t_1, \ldots, t_l\}$
be a set of generators of $M/N$, so that $M/N \cong A^l/R$, where $R$ is the
submodule of relations among the $t_j$; let $\{(a_{i1}, \ldots, a_{il}) \mid 1 \le i \le s\}$ be a
set of generators of $R$. Choose a pre-image $\bar{t}_j$ of $t_j$ in $M$ for each $j$. Then
set $\eta_i = \sum_j a_{ij}\bar{t}_j \in N$ (for $1 \le i \le s$). By the preceding theorem there exist
$\zeta_j \in N$ such that $\eta_i = \sum_j a_{ij}\zeta_j$ (for $1 \le i \le s$). Then $\sum_j a_{ij}(\bar{t}_j - \zeta_j) = 0$ (for
$1 \le i \le s$), and setting $f(t_j) = \bar{t}_j - \zeta_j$, we obtain a linear map $f\colon M/N \to M$
which satisfies the requirement.
\end{proof}
